\documentclass{beamer}

\usepackage{pgfplots}
\usepackage{hyperref}

\title[Determinants]{Determinants}
\subtitle{MATH 1554: \textit{Linear Algebra}}
\author{Ahmed Mohammed}
\date{PLUS Session, June 2023}
\institute{Georgia Institute of Technology}

\usetheme{Berlin}
\AtBeginSection[] {
	\begin{frame}
		\frametitle{Table Contents}
		\tableofcontents[currentsection]
	\end{frame}
}

\DeclareMathOperator{\vspan}{Span}

\begin{document}
\frame{\titlepage}

\section{About Subspaces of \(\mathbb{R}^n\)}
\begin{frame}
	\frametitle{Opening Activity: Revision}
	Imagine I have two vectors linearly independent vectors, \(\vec{v}, \vec{u} \in \mathbb{R}^3\).\pause 

	Notice that they span a plane. \pause
	\begin{figure}[!h]
		\begin{tikzpicture}
			\begin{axis}[no markers, xticklabels={}, yticklabels={}, zticklabels={}, width=6cm, height=4cm, xlabel = \(x_1\), ylabel = \(x_2\), zlabel=\(x_3\)]
				\addplot3[surf, domain=-10:10, color=blue!60, shader=flat]{-3*x + 2*y};
			\end{axis}
		\end{tikzpicture}
	\end{figure}\pause
	Is it true that \[\vspan\{\vec{v}, \vec{u}\} = \mathbb{R}^2?\]
\end{frame}

\begin{frame}
	\frametitle{Opening Activity: Revision (\textit{Cont.})}
	No!
	Two ways to think about it: \pause
		\begin{itemize}
			\item<1-> Algebraically
				\begin{itemize}
					\item<3-> \(\vspan \{\vec{v}, \vec{u}\} = \{c_1\vec{v}_1 + c_2\vec{v}_2 \mid c_1, c_2 \in \mathbb{R}\} \subset \mathbb{R}^3 \ne \mathbb{R}^2\).
				\end{itemize}
			\item<2-> ``Geometrically"
				\begin{itemize}
					\item<4-> A plane in \(\mathbb{R}^3\) is a collection of vectors whose tip lie on the plane. Since the vectors are in 
							\(\mathbb{R}^3\), a plane in \(\mathbb{R}^3\) is a collection of vectors in \(\mathbb{R}^3\).
				\end{itemize}
		\end{itemize} \pause

		\onslide<5->{
			Main confusion seems to stem from the fact that \(\mathbb{R}^2\) is a ``plane", and so is a plane in \(\mathbb{R}^3\). Hence they must be
			the same.
		}

		\onslide<6->{
			\begin{alertblock}{``Structure" vs. Content}
				Different subspaces \(U\) and \(V\) of different vector spaces can be ``similar in overall structure" but different in content.
			\end{alertblock}
		}
\end{frame}

\begin{frame}
	\begin{alertblock}{``Structure" vs. Content}
		Different subspaces \(U\) and \(V\) of different vector spaces can be ``similar in overall structure" but different in content.
	\end{alertblock}
	Examples include
		\begin{itemize}
			\item<2-> \(\mathbb{R}\) vs.\! a line in \(\mathbb{R}^n,\) \(n > 1\).
			\item<3-> \(\mathbb{R}^2\) vs.\! a plane in \(\mathbb{R}^3\).
			\item<4-> \(\mathbb{R}^3\) vs.\! a three-dimensional subspace of \(\mathbb{R}^4\).
		\end{itemize}
	\onslide<5->{Notice that each pair has the same dimension. Implies connection between dimension of two subspaces and how "similar" they look.}
	
	\onslide<6->{
		This property of ``similarity" is nice because if you're trying to solve a question in \(U\), and \(U\) is similar to \(V\), then you can find an equivalent question in \(V\).
	}
\end{frame}


\section{\href{https://www.geogebra.org/m/v9dyypp9}{Coordinates}}
\begin{frame}
	\frametitle{Coordinates and ``Similarity"}
	Going back to the example of the plane in \(\mathbb{R}^3\), notice that if we picked a basis for the subspace, we can identify each vector in the subspace by its 
	coordinates. That is, since for any \(\vec{x} \in \mathbb{R}^3\) that's in the plane, 
	\[\vec{x} = c_ 1\vec{b}_1 + c_2\vec{b}_2 = c_1\begin{bmatrix}1 \\ 2 \\ 1\end{bmatrix} + c_2\begin{bmatrix}0 \\ 1 \\ 2\end{bmatrix},\]
	we can \emph{uniquely} idenitfy \(\vec{x}\) with the scalars \(c_1, c_2\), which we'll arrange as a vector \([\vec{x}]_{\mathcal{B}} =\begin{bmatrix}c_1 \\ c_2\end{bmatrix}.\) 

	We've implicitly shown that the plane in \(\mathbb{R}^3\) and \(\mathbb{R}^2\) are ``similar." 
\end{frame}

\begin{frame}
\frametitle{Why ``Similarity" Is Useful}
We can ``switch" between \(\mathbb{R}^2\) and the plane subspace in \(\mathbb{R}^3\) in multiple ways, including:
\begin{itemize}
\item If \(\{\vec{x}_1, \vec{x}_2, \dots, \vec{x}_p\}\) is linearly in/dependent, so is \(\{[\vec{x}_1]_{\mathcal{B}}, [\vec{x}_2]_{\mathcal{B}}, \dots, [\vec{x}_p]_{\mathcal{B}}\}\).
	Vice versa.
\item \(\vec{x} + \vec{y} = \vec{z} \iff [\vec{x}]_{\mathcal{B}} + [\vec{y}]_{\mathcal{B}} = [\vec{z}]_{\mathcal{B}}.\)
\end{itemize}
\end{frame}

\section{Determinants}
\begin{frame}
	\frametitle{The Whole Point of This Session}
	The determinant is a multilinear function. It's linear ``on one argument." It takes a square matrix and spits out a number. Some properties it posseses is: \pause
	\begin{itemize}
		\item \(\det A = -\det B\) for any \(A, B \in \mathbb{R}^{n \times n}\), where \(A, B\) differ by a row swap. \pause
		\item Scaling one of the rows of \(A\) by \(c \in \mathbb{R}\) scales the determinant by that amount. \pause
		\item Adding multiple of row to another row doesn't alter the determinant. \pause
		\item \(\det(AB) = \det A \cdot \det B.\) (Inluding elementary matrices.) \pause
		\item \(\det A = \det A^T\).
	\end{itemize}
\end{frame}
 
\begin{frame}
	\frametitle{Closing Activity: Conceptual Questions About Determinants}
	\begin{itemize}
		\item If \(\det A = \begin{vmatrix}a & b & c \\ d & e & f \\ g & h & i\end{vmatrix} = 5,\) then what's \(\begin{vmatrix}a & a - c & c \\ d & d - f & f \\ g & g- i & i\end{vmatrix}\)? \pause 
		\item Consider \(A, B \in \mathbb{R}^{n \times n}\) that are row equivalent. The rows of \(A\) are made by swaping a row of \(B\) and scaling another row by a factor of 
			2. If \(\det A = 5\), what is \(\det B\)? \pause
		\item If \(A = LU \in \mathbb{R}^{n \times n}\), and the determinant of \(U\) is 6, what's the determinant of \(A\)? \pause
		\item If a matrix \(A \in \mathbb{R}^{n \times n}\) has been row-reduced to the corresponding identity matrix by \emph{only} adding multiples of a row to the one below,
			what's the determinant of \(A\)? \pause
	\end{itemize}
\end{frame}
\end{document}